\documentclass[../../main.tex]{subfiles}
\graphicspath{{../../image/}{../../figure/}}
\begin{document}

\section{Cauchy 问题：经典解的存续与消亡}\label{section:ch06-cauchy问题经典解}

我们考虑齐次标量守恒律的 Cauchy 问题：
\begin{align}
\partial_{t}u(x,t)+\operatorname{div}G(u(x,t))=0,\qquad x\in\mathbb{R}^{m},\ t>0,
\label{align:ch06-6-1-1}
\end{align}
\begin{align}
u(x,0)=u_{0}(x),\qquad x\in\mathbb{R}^{m}.
\label{align:ch06-6-1-2}
\end{align}
流通量 $G(u)=(G_{1}(u),\cdots,G_{m}(u))$ 是定义在 $\mathbb{R}$ 上的给定光滑函数，取值于 $\mathbb{R}^{1\times m}$。

与\eqref{align:ch06-6-1-1}的连续可微解 $u$ 相联系的特征线，是如下常微分方程组的一条轨道 $\xi\colon[0,T)\to\mathbb{R}^{m}$：
\begin{align}
\frac{\mathrm{d}x}{\mathrm{d}t}=G^{\prime}(u(x,t))^{\top}.
\label{align:ch06-6-1-3}
\end{align}
对每条特征线 $\xi$，我们关联微分算子
\begin{align}
\frac{\mathrm{d}}{\mathrm{d}t}=\partial_{t}+G^{\prime}(u(\xi(t),t))\operatorname{grad},
\label{align:ch06-6-1-4}
\end{align}
它确定沿 $\xi$ 的方向导数。特别地，由于 $u$ 满足\eqref{align:ch06-6-1-1}，$\mathrm{d}u/\mathrm{d}t=0$，即 $u$ 沿任一特征线为常数。由\eqref{align:ch06-6-1-3}可知，特征线的斜率为常数。于是所有特征线都是直线，而解沿其保持常值。借助这一性质，可对\eqref{align:ch06-6-1-1}、\eqref{align:ch06-6-1-2}的经典解作细致研究。特别地，对标量守恒律，\cref{theorem:ch05-凸熵局部存在}容许如下精细化：


\begin{theorem}\label{theorem:ch06-6-1-1}
设定义在 $\mathbb{R}^{m}$ 上的 $u_{0}$ 有界且 Lipschitz 连续。令
\begin{align}
\kappa=\operatorname*{ess\,inf}_{y\in\mathbb{R}^{m}}\operatorname{div}G^{\prime}(u_{0}(y)).
\label{align:ch06-6-1-5}
\end{align}
则在最大区间 $[0,T_{\infty})$ 上存在\eqref{align:ch06-6-1-1}、\eqref{align:ch06-6-1-2}的经典解 $u$，其中当 $\kappa\geqslant 0$ 时 $T_{\infty}=\infty$，当 $\kappa<0$ 时 $T_{\infty}=-\kappa^{-1}$。此外，若 $u_{0}$ 为 $C^{k}$，则 $u$ 亦然。
\end{theorem}

\begin{proof}
先设 $u_{0}\in C_{0}^{\infty}(\mathbb{R}^{m})$，并令 $u$ 为\eqref{align:ch06-6-1-1}、\eqref{align:ch06-6-1-2}的唯一光滑经典解，按\cref{theorem:ch05-凸熵局部存在}在最大时间区间 $[0,T_{\infty})$ 上定义。

由上述特征线的性质，对任一点 $(x,t)\in\mathbb{R}^{m}\times[0,T_{\infty})$，存在 $y\in\mathbb{R}^{m}$ 使得
\begin{align}
y=x-tG^{\prime}(u(x,t))^{\top},\qquad u(x,t)=u_{0}(y).
\label{align:ch06-6-1-6}
\end{align}
由\eqref{align:ch06-6-1-6}易得
\begin{align}
\nabla u(x,t)=\frac{\nabla u_{0}(y)}{1+t\operatorname{div}G^{\prime}(u_{0}(y))},\qquad
\partial_{t}u(x,t)=\frac{-G^{\prime}(u_{0}(y))\cdot\nabla u_{0}(y)}{1+t\operatorname{div}G^{\prime}(u_{0}(y))},
\label{align:ch06-6-1-7}
\end{align}
特别地，这蕴含：若 $\kappa\geqslant 0$ 则 $T_{\infty}=\infty$，若 $\kappa<0$ 则 $T_{\infty}=-\kappa^{-1}$。

现设 $u_{0}$ 仅在 $\mathbb{R}^{m}$ 上 Lipschitz 连续。令当 $\kappa\geqslant 0$ 时 $T_{\infty}=\infty$，当 $\kappa<0$ 时 $T_{\infty}=-\kappa^{-1}$。借助磨光子，构造 $C_{0}^{\infty}(\mathbb{R}^{m})$ 中的序列 $\{u_{0n}\}$，使其在紧集上一致收敛于 $u_{0}$，且在 $\nabla u_{0}$ 的每个 Lebesgue 点 $y$ 处 $\nabla u_{0n}(y)$ 趋于 $\nabla u_{0}(y)$。

以 $u_{0n}$ 为初值的\eqref{align:ch06-6-1-1}的经典解 $u_{n}$ 定义于最大时间区间 $[0,T_{n})$。令 $T=\liminf_{n\to\infty}T_{n}$，注意 $0<T\leqslant T_{n}$。由\eqref{align:ch06-6-1-7}，$u_{n}$ 在 $\mathbb{R}^{m}\times[0,T)$ 的每个紧子集上等 Lipschitz，因而 $\{u_{n}\}$ 的某子列在紧集上一致收敛于局部 Lipschitz 函数 $u$。显然 $u$ 从 $\{u_{n}\}$ 继承性质\eqref{align:ch06-6-1-6}。这进而蕴含：若 $u_{0}$ 在某点 $y$ 可微，则 $u$ 沿特征线 $x=y+tG^{\prime}(u_{0}(y))$ 可微，且导数由\eqref{align:ch06-6-1-7}给出。特别地，$u$ 是\eqref{align:ch06-6-1-1}、\eqref{align:ch06-6-1-2}在时间区间 $[0,T)$ 上的经典解。若 $T=T_{\infty}$，\eqref{align:ch06-6-1-7}蕴含 $[0,T_{\infty})$ 是最大时间区间，定理结论得证。反之，若 $T<T_{\infty}$，\eqref{align:ch06-6-1-7}蕴含 $u$ 可延拓至 $t=T$，且 $u(\cdot,T)$ 在 $\mathbb{R}^{m}$ 上 Lipschitz。于是可重复该过程，将 $u$ 的存在时间区间延长至 $[0,T_{\infty})$，而它必然是最大的。

最后，对\eqref{align:ch06-6-1-6}应用隐函数定理，可知当 $u_{0}$ 为 $C^{k}$ 时解 $u$ 亦为 $C^{k}$。证明完成。
\end{proof}

由上述讨论可知，经典解的寿命一般是有限的。因而必须处理弱解。

另一种富有启发性地考察\eqref{align:ch06-6-1-1}经典解 $u$ 的方式，是将其实现为函数 $f(v;x,t)$ 的“等值面”；$f$ 定义在 $\mathbb{R}\times\mathbb{R}^{m}\times\mathbb{R}$ 上，即当 $u$ 满足\eqref{align:ch06-6-1-1}时
\begin{align}
f(u(x,t);x,t)=0.
\label{align:ch06-6-1-8}
\end{align}
容易看出，为此 $f$ 必须满足输运方程
\begin{align}
\partial_{t}f(v;x,t)+\sum_{\alpha=1}^{m}G_{\alpha}^{\prime}(v)\partial_{\alpha}f(v;x,t)=0.
\label{align:ch06-6-1-9}
\end{align}
于是我们把非线性方程\eqref{align:ch06-6-1-1}化为线性方程，代价是自变量个数从 $m+1$ 增至 $m+2$。特别地，为求解初值问题\eqref{align:ch06-6-1-1}、\eqref{align:ch06-6-1-2}，应求解\eqref{align:ch06-6-1-9}带初始条件 $f(v;x,0)=v-u_{0}(x)$ 的 Cauchy 问题。由于\eqref{align:ch06-6-1-9}是线性的，该 Cauchy 问题在 $\mathbb{R}\times\mathbb{R}^{m}\times\mathbb{R}$ 上存在解。所得 $f$ 将通过\eqref{align:ch06-6-1-8}诱导出\eqref{align:ch06-6-1-1}、\eqref{align:ch06-6-1-2}的经典解 $u$，其有效范围直至 $f_{v}$ 首次为零。我们将在\cref{section:ch06-动力论表述}中、于弱解的框架下回到输运方程\eqref{align:ch06-6-1-9}。


\end{document}
